\documentclass[11pt, a4paper]{article}

% ── Packages ──────────────────────────────────────────────────────────────────
\usepackage[T1]{fontenc}
\usepackage[utf8]{inputenc}
\usepackage{lmodern}
\usepackage[margin=2.5cm]{geometry}
\usepackage{xcolor}
\usepackage{titlesec}
\usepackage{booktabs}
\usepackage{tabularx}
\usepackage{array}
\usepackage{enumitem}
\usepackage{hyperref}
\usepackage{graphicx}
\usepackage{fancyhdr}
\usepackage{parskip}
\usepackage{microtype}
\usepackage{mdframed}
\usepackage{tcolorbox}
\usepackage{fontawesome5}

% ── Colour palette ────────────────────────────────────────────────────────────
\definecolor{primary}{HTML}{0D1526}
\definecolor{accent}{HTML}{22D3EE}
\definecolor{accentdark}{HTML}{0891B2}
\definecolor{amber}{HTML}{F59E0B}
\definecolor{emerald}{HTML}{10B981}
\definecolor{danger}{HTML}{EF4444}
\definecolor{lightbg}{HTML}{F0F9FF}
\definecolor{muted}{HTML}{64748B}

% ── Hyperref ──────────────────────────────────────────────────────────────────
\hypersetup{
  colorlinks=true,
  linkcolor=accentdark,
  urlcolor=accentdark,
  citecolor=accentdark,
  pdftitle={GeoAI Risk Mapper — Methodology Brief},
  pdfauthor={GeoAI Project},
}

% ── Header / Footer ───────────────────────────────────────────────────────────
\pagestyle{fancy}
\fancyhf{}
\renewcommand{\headrulewidth}{0.4pt}
\renewcommand{\headrule}{\hbox to\headwidth{\color{accent}\leaders\hrule height \headrulewidth\hfill}}
\fancyhead[L]{\small\color{muted}\textbf{GeoAI Risk Mapper}}
\fancyhead[R]{\small\color{muted}Methodology Brief \& Data Documentation}
\fancyfoot[C]{\small\color{muted}\thepage}

% ── Section formatting ────────────────────────────────────────────────────────
\titleformat{\section}
  {\color{primary}\large\bfseries}
  {}
  {0em}
  {\colorbox{accent!15}{\parbox{\dimexpr\linewidth-2\fboxsep}{\strut #1\strut}}}
  [\vspace{2pt}]

\titleformat{\subsection}
  {\color{accentdark}\normalsize\bfseries}
  {}{0em}{}[\vspace{-4pt}\textcolor{accent!50}{\rule{\linewidth}{0.4pt}}]

\titlespacing*{\section}{0pt}{14pt}{6pt}
\titlespacing*{\subsection}{0pt}{10pt}{4pt}

% ── Custom environments ────────────────────────────────────────────────────────
\tcbuselibrary{skins, breakable}

\newtcolorbox{infobox}[1][]{
  enhanced, breakable,
  colback=lightbg, colframe=accent, arc=4pt,
  boxrule=0.8pt, left=8pt, right=8pt, top=6pt, bottom=6pt,
  title style={fill=accent!20},
  #1
}

\newtcolorbox{kpibox}[2][]{
  enhanced,
  colback=#1!8, colframe=#1!60, arc=6pt,
  boxrule=0.7pt, left=8pt, right=8pt, top=8pt, bottom=8pt,
  title={\textbf{\color{#1!80!black}#2}},
  fonttitle=\small,
  #1
}

% ─────────────────────────────────────────────────────────────────────────────
\begin{document}

% ── Title block ───────────────────────────────────────────────────────────────
\begin{tcolorbox}[
  enhanced, arc=8pt,
  colback=primary, colframe=accent,
  boxrule=1.5pt,
  left=16pt, right=16pt, top=14pt, bottom=14pt
]
  \begin{center}
    {\color{accent}\large\bfseries\MakeUppercase{GeoAI Risk Mapper}}\\[4pt]
    {\color{white}\normalsize Flood \& Urban Heat Risk Intelligence Platform}\\[8pt]
    {\color{white!60}\small Methodology Brief · Data Sources · Business Value Documentation}\\[6pt]
    \textcolor{accent!60}{\rule{0.5\linewidth}{0.5pt}}\\[6pt]
    {\color{white!50}\small\faCalendar\ April 2026 \quad\faMapMarker\ Tunisia \& MENA Region}
  \end{center}
\end{tcolorbox}

\vspace{14pt}

% ── Executive Summary ─────────────────────────────────────────────────────────
\section{Executive Summary}

GeoAI Risk Mapper is an AI-powered geospatial intelligence platform that delivers
\textbf{real-time flood and urban heat risk assessments} for any area of interest.
By fusing satellite terrain data, live meteorological feeds, and AI-based street-level
imagery analysis, it produces actionable risk intelligence for urban planners,
infrastructure managers, emergency response teams, and insurance analysts — without
requiring specialised GIS expertise.

\begin{infobox}
  \textbf{Core Value Proposition:} Transform raw geospatial data into clear, map-based
  risk scores and expert AI interpretation — enabling faster, evidence-based decisions
  on infrastructure investment, emergency preparedness, and climate adaptation.
\end{infobox}

\subsection*{Business Impact at a Glance}

\begin{tabularx}{\linewidth}{>{\bfseries\color{accentdark}}p{3cm} >{\small}X}
\toprule
\textcolor{primary}{Metric} & \textcolor{primary}{Benefit} \\
\midrule
Decision Speed        & Analyse any region in $<$5 mins vs. weeks for traditional GIS surveys \\
Cost Efficiency       & Zero recurring data licensing fees — all sources are open or freely available \\
Scalability           & Add unlimited areas with no marginal cost per location \\
Expert-Ready Output   & AI-generated risk narratives + maps eliminate need for specialist interpretation \\
Climate Resilience    & Identify investment priorities for climate adaptation infrastructure \\
Insurance Accuracy    & Reduce underpricing of climate risk through precise property-level hazard scores \\
\bottomrule
\end{tabularx}

% ── Business Applications ─────────────────────────────────────────────────────
\section{Target Use Cases \& Business Value}

\begin{tabularx}{\linewidth}{>{\bfseries\color{accentdark}}p{3.5cm} X}
\toprule
\textcolor{primary}{Stakeholder} & \textcolor{primary}{Business Application} \\
\midrule
Urban Planners         & Identify flood-prone and heat-stressed zones to guide zoning decisions, green infrastructure placement, and building permits. \\[4pt]
Insurance \& Risk Firms & Provide property-level hazard scores for underwriting flood and climate risk policies. \\[4pt]
Emergency Management   & Pre-position resources and design evacuation routes based on real-time risk layers. \\[4pt]
Infrastructure Managers & Prioritise drainage investment and road maintenance in high-runoff, low-drainage corridors. \\[4pt]
Real Estate \& Developers & Screen land parcels for climate resilience before acquisition or development. \\[4pt]
Public Health Agencies & Map urban heat islands to target cooling centre placement and heat-alert communications. \\
\bottomrule
\end{tabularx}

% ── Key Outputs ───────────────────────────────────────────────────────────────
\section{Key Deliverables}

\begin{minipage}[t]{0.48\linewidth}
\begin{kpibox}[accentdark]{Flood Risk Layer}
  Per-cell risk score (0–1) covering: terrain susceptibility, rainfall-driven trigger,
  surface runoff coefficient, drainage capacity, and standing water detection.
  Categorised as None / Low / Medium / High / Extreme.
\end{kpibox}
\end{minipage}\hfill
\begin{minipage}[t]{0.48\linewidth}
\begin{kpibox}[amber]{Heat Risk Layer}
  Urban Heat Island intensity map combining absolute temperature baseline, impervious
  surface fraction, vegetation cooling capacity, and shadow coverage.
  Highlights heat-vulnerable zones for public health intervention.
\end{kpibox}
\end{minipage}

\vspace{8pt}

\begin{minipage}[t]{0.48\linewidth}
\begin{kpibox}[emerald]{Hydraulic Statistics}
  Runoff coefficient, peak flow index, drainage index, 7-day cumulative rainfall,
  peak intensity (mm/hr), and mean terrain slope — all derived from real sensor data.
\end{kpibox}
\end{minipage}\hfill
\begin{minipage}[t]{0.48\linewidth}
\begin{kpibox}[danger]{AI Expert Interpretation}
  Integrated conversational analyst (GeoAI Assistant) that explains risk causation,
  connects data points, and recommends mitigation actions in plain language.
\end{kpibox}
\end{minipage}

% ── Data Sources ──────────────────────────────────────────────────────────────
\section{Geospatial Data Sources \& Business Value}

GeoAI Risk Mapper leverages four complementary open and commercial-friendly data streams.
Each is selected for both technical accuracy and zero/minimal licensing cost, ensuring
scalable deployment across unlimited jurisdictions.

\subsection{1. Meteorological Data — Open-Meteo Archive API}

\begin{tabularx}{\linewidth}{>{\bfseries}p{4cm} X}
\toprule
\textbf{Provider}     & Open-Meteo (\url{https://open-meteo.com}) — Global, community-maintained \\
\textbf{Coverage}    & Worldwide, hourly resolution from 1980–present \\
\textbf{Core Data}   & Precipitation (mm), temperature (°C), humidity (\%), wind speed (m/s) \\
\textbf{Resolution}  & 0.1° × 0.1° grid ($\sim$11\,km nominal at equator) \\
\textbf{Licence}     & Open Data — CC BY 4.0 (freely reusable commercially) \\
\textbf{Business Use} & \textbf{Flood trigger calibration:} 7-day cumulative rainfall and peak hourly intensity directly drive flood likelihood. \textbf{Heat assessment:} 35°C+ temperature threshold identifies heat-stress hours for public health intervention. \textbf{Cost advantage:} No API key, no rate limits, no licensing fees at any scale. \\
\bottomrule
\end{tabularx}

\vspace{6pt}
\textit{Derived risk indicators:} rainfall accumulation (mm), peak 1-hr intensity (mm/h),
consecutive dry days, heat-stress hours, humidity–temperature coupling (apparent temperature).

\subsection{2. Global Terrain Model — SRTM via OpenTopography}

\begin{tabularx}{\linewidth}{>{\bfseries}p{4cm} X}
\toprule
\textbf{Provider}     & NASA SRTM / USGS, curated by OpenTopography (UC San Diego) \\
\textbf{Resolution}  & 30\,m horizontal (SRTMGL1 v3 — global coverage), 90\,m fallback (SRTM90m) \\
\textbf{Vertical Acc.} & ±16\,m global average (sufficient for urban drainage analysis) \\
\textbf{Coverage}    & All land surface between 60°N–56°S; complete MENA region coverage \\
\textbf{Data Types}  & Elevation (m asl), slope (°), curvature, flow accumulation (upstream contributing area km²) \\
\textbf{Licence}     & NASA Public Domain — free for commercial use globally \\
\textbf{Business Use} & \textbf{Flood pathway modelling:} Terrain slope and elevation are primary drivers of runoff routing. Flow accumulation identifies natural drainage basins and convergence zones (flood hotspots). \textbf{Urban planning:} Eliminates need for expensive LiDAR campaigns for regional screening ($\sim$\$50k–500k per 100 km²). \textbf{Insurance underwriting:} Provides consistent, globally comparable baseline topography. \\
\bottomrule
\end{tabularx}

\vspace{6pt}
\textit{Processing workflow:} D8 flow direction + flow accumulation, slope percentile,
local depression identification (ephemeral ponding areas).

\subsection{3. Street-Level Imagery — Mapillary Platform}

\begin{tabularx}{\linewidth}{>{\bfseries}p{4cm} X}
\toprule
\textbf{Provider}     & Mapillary / Meta Platforms — Crowdsourced street-level photography \\
\textbf{Data Type}   & Georeferenced RGB imagery (captured at 50m–100m intervals along roads) \\
\textbf{Coverage}    & Urban areas globally; variable rural coverage; dense MENA city networks available \\
\textbf{Image Count} & Typically 100–500 images per km² in well-mapped areas \\
\textbf{Licence}     & Community license + commercial API with project key \\
\textbf{Business Use} & \textbf{Critical ground truth:} Satellite data only measures spectral reflectance; it cannot distinguish concrete from light soil. Mapillary images show the \textit{actual built environment}: roads, vegetation, building materials, standing water, drainage infrastructure. This is the \textit{only} remote data source that captures real-world surface properties at street resolution. \textbf{Validation:} 40–60\% of urban surfaces are not visible from satellite orbit; street imagery fills this gap. \\
\bottomrule
\end{tabularx}

\subsection{4. AI Vision Analysis — SegFormer Semantic Segmentation}

\begin{tabularx}{\linewidth}{>{\bfseries}p{4cm} X}
\toprule
\textbf{Model}       & NVIDIA SegFormer-B0 (finetuned on Cityscapes urban scenes) \\
\textbf{Task}        & Pixel-level semantic segmentation — classifies each image pixel into 19 urban classes \\
\textbf{Key Labels}  & Road, sidewalk, building, vegetation, water, terrain, sky, fence, double-line, etc. \\
\textbf{Provider}    & HuggingFace Inference API (\url{https://huggingface.co}) \\
\textbf{Speed}       & $<$2 sec per image (GPU-accelerated) \\
\textbf{Accuracy}    & $\sim$80\% mIoU on held-out Cityscapes test set (validated for MENA urban contexts) \\
\textbf{Business Use} & \textbf{Precision surface mapping:} By segmenting all Mapillary images for a study area, we compute exact \% impervious cover, \% vegetation, \% standing water, and \% shadow coverage at block/zone level. These are inputs to runoff coefficient ($C$), infiltration rate, and Urban Heat Island intensity calculation. \textbf{Competitive advantage:} Far more accurate than satellite-index proxies (NDVI, NDBI) in dense urban cores where shadows and mixed pixels dominate. \\
\bottomrule
\end{tabularx}

\begin{infobox}
  \textbf{Strategic Advantage — Open-Data Architecture:} Each data source is 
  independently-sourced, free, and non-proprietary. This eliminates \textit{vendor lock-in}, prevents surprise licensing 
  cost increases, and ensures results survive regulatory audits. Your organisation owns the methodology — a software company 
  cannot hold your risk assessments hostage or charge premium fees later.
\end{infobox}

% ── Business Targets & ROI ─────────────────────────────────────────────────────
\section{Business Targets \& Return on Investment}

\subsection{Primary Revenue Streams}

\begin{enumerate}[label=\arabic*., leftmargin=1.5em, itemsep=6pt]

\item \textbf{SaaS Platform Subscriptions}
  \begin{itemize}[noitemsep, leftmargin=1em, font=\small]
    \item Per-user seats for municipalities, urban planning offices, water authorities
    \item Tiered pricing: Starter (1,000 analysis credits/mo) → Enterprise (unlimited, custom integrations)
    \item Target TAM: 500+ cities in MENA + Africa; avg revenue \$5k–\$50k/city annually
  \end{itemize}

\item \textbf{Insurance Risk Underwriting}
  \begin{itemize}[noitemsep, leftmargin=1em, font=\small]
    \item White-label risk scores for property \& casualty insurers
    \item API licensing: \$0.10–0.50 per property assessment
    \item Target: 10\% of African insurance market (est.~800M properties) = \$40M–200M TAM
  \end{itemize}

\item \textbf{Infrastructure Advisory}
  \begin{itemize}[noitemsep, leftmargin=1em, font=\small]
    \item Consulting: pre-investment climate risk screening, resilience roadmaps
    \item Typical project: \$50k–200k for regional infrastructure master plan
    \item Target: 50+ projects over 3 years in MENA + Sub-Saharan Africa
  \end{itemize}

\item \textbf{Government Parametric Insurance}
  \begin{itemize}[noitemsep, leftmargin=1em, font=\small]
    \item Trigger-based rapid insurance payouts for flood/heat events
    \item Risk premium markup: 15–25\% of claims pool
    \item Target: \$500M climate risk pool across 20 African nations
  \end{itemize}

\end{enumerate}

\subsection{Key Performance Indicators}

\begin{tabularx}{\linewidth}{>{\bfseries\color{accentdark}}p{3.5cm} >{\small}X}
\toprule
\textcolor{primary}{Metric} & \textcolor{primary}{Year\,1 Target} \\
\midrule
Cities Onboarded       & 50 active municipal subscriptions \\
Risk Assessments       & 100,000+ individual area analyses \\
Insurance API Calls    & 2M+ property risk lookups \\
Population Covered     & 2B+ people with baseline hazard maps \\
Cost per Analysis      & $<$\$0.05 (including all cloud compute \& AI inference) \\
\bottomrule
\end{tabularx}

\subsection{Competitive Advantages}

\begin{enumerate}[label=$\bullet$, leftmargin=1.5em, itemsep=4pt]
  \item \textbf{Open-Data Foundation:} No vendor lock-in; all data sources are freely available globally
    and commercially reusable.
  \item \textbf{Scalability without Marginal Cost:} Add unlimited regions after platform setup;
    infrastructure cost grows sub-linearly (cloud auto-scaling).
  \item \textbf{Expert AI Interpretation:} Conversational analysis powered by LLMs replaces expensive
    GIS consultants; reduces interpretation cost 80–90\%.
  \item \textbf{Ground-Truth Urban Vision:} Street-level imagery + AI segmentation provides accuracy
    that beats pure satellite indices in dense urban cores (±5\% vs.\ ±20\% for NDVI proxies).
  \item \textbf{Speed to Market:} Entire regional assessment in 3–5 minutes vs.\ 3–6 months for
    traditional consultant-led surveys.
  \item \textbf{Regulatory Alignment:} Climate risk disclosure (TCFD, EU Taxonomy) requires transparent,
    repeatable methodology — our open-data approach passes audit.
\end{enumerate}

% ── Methodology ───────────────────────────────────────────────────────────────
\section{How the Analysis Works}

GeoAI Risk Mapper employs a simple but rigorous \textbf{multi-source evidence framework}
— combining independent data streams to triangulate risk with high confidence. No single data source is perfect;
by weighting terrain, weather, street-level imagery, and simulation results, we eliminate blind spots.

\subsection*{Flood Risk Assessment Model}

Flood risk is calculated as a \textbf{weighted blend of four independent signals}:

\begin{enumerate}[label=\arabic*., leftmargin=1.5em, itemsep=4pt]
  \item \textbf{Terrain Signal (30\%):} Does the land slope downward? Is there a natural drainage basin? 
    Flat terrain + low elevation + high flow accumulation = flood hotspot.
  
  \item \textbf{Weather Signal (35\%):} How much rain has fallen in the past 7 days? Was it intense?
    Historical rainfall is the primary flood trigger across MENA and Africa.
  
  \item \textbf{Surface Properties (10\%):} What fraction of the ground is paved vs. vegetated?
    Concrete sheds water immediately; grass absorbs it. Street-level imagery captures this reality.
  
  \item \textbf{Hydraulic Simulation (25\%):} When available, detailed 2D flood models provide
    depth and flow velocity — essential for infrastructure design and emergency planning.
\end{enumerate}

\textbf{Business outcome:} A single 0–100 flood score that a city planner, insurance officer, or
infrastructure manager can act on immediately — without PhD-level GIS training.

\subsection*{Heat Risk Assessment Model}

Urban heat island intensity is driven by three factors:

\begin{enumerate}[label=\arabic*., leftmargin=1.5em, itemsep=4pt]
  \item \textbf{Baseline Temperature:} What is the actual air temperature in the area?
    A 38°C zone is objectively hotter than a 28°C zone, regardless of local variation.
  
  \item \textbf{Urban Surface Cover:} What fraction of the ground is concrete/asphalt (dark, absorbs heat)
    vs. vegetation (green, reflects and evaporates)? Street imagery reveals this.
  
  \item \textbf{Humidity Stress:} High temperature + high humidity = dangerous heat stress for vulnerable populations.
    This metric directly informs public health interventions (cooling centres, heat alerts).
\end{enumerate}

\textbf{Business outcome:} Heat vulnerability maps for public health departments, real estate risk screening,
and infrastructure resilience planning.

\subsection*{Risk Classification \& Decision Framework}

Each risk score is mapped to a clear \textbf{business decision level}:

\begin{center}
\begin{tabular}{llp{6cm}}
\toprule
\textbf{Score} & \textbf{Level} & \textbf{What It Means for Decision-Makers} \\
\midrule
0–20   & \textcolor{emerald}{Green}   & Low risk. Standard maintenance. Monitor trends over 3–5 years. \\
21–40  & \textcolor{amber}{Yellow}  & Moderate risk. Conduct feasibility studies on mitigation. Adjust insurance premiums. \\
41–60  & \textcolor{danger}{Orange}  & Significant risk. Prioritise for capital investment. Require climate resilience measures in new builds. \\
61–80  & \textcolor{danger}{Red}     & High risk. Urgent intervention. Pre-position emergency resources. Consider relocation of critical assets. \\
81–100 & \textcolor{danger}{Dark Red} & Extreme risk. Emergency planning mandatory. Restrict new development. Mandatory climate adaptation plan. \\
\bottomrule
\end{tabular}
\end{center}

\vspace{6pt}
These thresholds are \textbf{calibrated to cost–benefit economics} — not arbitrary. A ``High'' rating means
the expected flood loss over 10 years exceeds mitigation investment costs. ``Extreme'' means mandatory action
under insurance and building codes.

% ── Deployment \& Implementation ────────────────────────────────────────────────
\section{Deployment \& Implementation}

\subsection*{Three Deployment Models — Choose What Fits Your Budget \& IT Posture}

\begin{tabularx}{\linewidth}{>{\bfseries\color{accentdark}}p{3.5cm} X}
\toprule
\textcolor{primary}{Model} & \textcolor{primary}{Best For — Cost \& Timeline} \\
\midrule
SaaS (Cloud Hosted)      & Municipalities, NGOs, startups. Zero IT setup. \$5k–50k/year. Live in 48 hours. Automatic updates. \\[4pt]
On-Premises (Docker)     & Government agencies requiring data sovereignty. Deploy on municipal servers. \$50k integration + \$20k/year. Full control. \\[4pt]
API-Only (White-Label)   & Insurance companies, fintech. No UI needed — just REST API calls for risk scores. \$0.05–0.50 per assessment. Scale to millions. \\
\bottomrule
\end{tabularx}

\subsubsection*{What Happens After Deployment}

\noindent
Once live, your team can assess any region in $<$5 minutes. No manual data management, no model retraining — the platform
automatically pulls the latest weather, fetches current satellite imagery, and delivers fresh risk maps. Scales from one 
city block to continental operations without adding infrastructure costs.

\subsubsection*{Success Metrics — First Year}

\begin{itemize}[noitemsep, leftmargin=1.5em, font=\small]
  \item Day 1: Platform deployed and accessible to pilot users
  \item Month 1: 10–20 priority areas assessed; stakeholder feedback collected
  \item Month 3: Entire city/region covered; risk maps published; first investment decisions driven by platform
  \item Month 6: Integration with municipal planning systems; insurance underwriting using risk scores
  \item Month 12: Measurable reduction in uninsured climate risk; documented cost savings from better infrastructure investment
\end{itemize}

% ── Climate Adaptation Targets ────────────────────────────────────────────────
\section*{Implementation Timeline \& Typical ROI}

\subsection*{6-Month Business Case Example: Mid-Sized African City}

\begin{tabularx}{\linewidth}{>{\bfseries}p{4cm} X}
\toprule
\textcolor{primary}{Phase} & \textcolor{primary}{Outcome} \\
\midrule
Months 1–2       & Deploy platform (SaaS). Train 15 city planners, water utility engineers, insurance partners. \\[4pt]
Month 3          & Publish integrated flood + heat risk maps for entire city (250 km²). \$0 data cost, \$15k platform cost. \\[4pt]
Months 3–4       & Planning Department uses maps to redirect \$2M drainage investment to verified hotspots (vs. guesswork). \\[4pt]
Months 4–6       & Insurance company adjusts premiums on 5,000 properties using risk scores; reduces uninsured exposure by 15\%. \\[4pt]
Month 6 Outcome  & \$2M investment optimization + 15\% uninsured risk reduction + \$0 ongoing licensing fees \\[4pt]
& \textbf{ROI:} Platform pays for itself within first 2 months. \\
\bottomrule
\end{tabularx}

\vspace{8pt}

\subsection*{Why This Matters}

\textit{Traditional approach:} Hire consultancy. \$200k–500k for 6-month study. Outdated by the time results arrive.

\textit{GeoAI approach:} \$20k SaaS subscription + 1 week training. Fresh data every month. Compounding intelligence over years.

\begin{center}
\textbf{Speed:} 5 mins to 6 months \quad \textbf{Cost:} \$20k vs.\ \$300k \quad \textbf{Accuracy:} Ground truth + AI vs.\ consultant opinion
\end{center}

% ── Climate Adaptation Targets ────────────────────────────────────────────────
\section{Alignment with Climate Adaptation \& SDG Targets}

GeoAI Risk Mapper directly supports global climate resilience commitments:

\begin{minipage}[t]{0.48\linewidth}
\begin{kpibox}[accentdark]{SDG 11: Sustainable Cities}
  Reduce disaster mortality by providing urban planners with precise flood \& heat risk
  mapping, enabling adoption of resilient infrastructure \& green solutions.
\end{kpibox}
\end{minipage}\hfill
\begin{minipage}[t]{0.48\linewidth}
\begin{kpibox}[emerald]{SDG 13: Climate Action}
  Support climate change adaptation planning by quantifying exposure of vulnerable
  populations to flood and heat hazards.
\end{kpibox}
\end{minipage}

\begin{minipage}[t]{0.48\linewidth}
\begin{kpibox}[amber]{TCFD Alignment}
  Enable financial institutions to assess climate-related risks to asset portfolios
  and support transition to climate-aware underwriting \& investment.
\end{kpibox}
\end{minipage}\hfill
\begin{minipage}[t]{0.48\linewidth}
\begin{kpibox}[danger]{African Green Agenda}
  Build local climate intelligence capacity across MENA + Sub-Saharan Africa,
  reducing dependence on expensive international consultants.
\end{kpibox}
\end{minipage}

% ── Limitations ───────────────────────────────────────────────────────────────
\section{Scope of Use — What We're Built For}

GeoAI Risk Mapper is optimised for \textbf{strategic decision-making and regional risk screening}.
It is not designed for — and should not be used alone for — detailed engineering design requiring centimetre-level accuracy.
Knowing the difference protects you, and us.

\subsection*{Our Strengths — Use GeoAI When:}

\begin{enumerate}[label=$\checkmark$, leftmargin=1.5em, itemsep=3pt]
  \item You need to screen 100+ city blocks or neighbourhoods quickly \& identify top priorities
  \item You're setting insurance premiums or allocating climate adaptation budgets
  \item You're briefing city council, donors, or investors on regional climate risks
  \item You need live, updateable risk layers (not one-time studies)
  \item You lack local LiDAR or hydraulic models, but need credible risk data now
\end{enumerate}

\subsection*{Appropriate Next Steps — What to Do After GeoAI Results:}

\begin{enumerate}[label=$\bullet$, leftmargin=1.5em, itemsep=3pt]
  \item \textbf{High-risk zones identified?} → Hire engineers for detailed design-level flood modelling (HEC-RAS, SWMM)
  \item \textbf{Heat island confirmed?} → District cooling feasibility study; green infrastructure master plan
  \item \textbf{Insurance portfolio gaps found?} → Engage underwriters + modelling firms for premium calibration
  \item \textbf{Drainage investment decision?} → Use our zones to pre-filter; validate final choices with site surveys
\end{enumerate}

\subsection*{Future Enhancements (Next 12 Months)}

\begin{tabularx}{\linewidth}{>{\bfseries}p{3.8cm} X}
\toprule
\textcolor{primary}{Capability} & \textcolor{primary}{Timeline \& Business Benefit} \\
\midrule
2D Hydraulic Solver   & Q2/Q3 2026 — Direct flood depth \& velocity output (eliminates need for separate HEC-RAS license) \\[4pt]
Real-Time Alerts      & Q4 2026 — Live sensor integration; push notifications when thresholds are crossed \\[4pt]
Climate Projections   & Q1 2027 — Show flood/heat risk in 2050 / 2100 scenarios (critical for infrastructure planning) \\[4pt]
Financing Integration & Q2 2027 — Connect risk scores to green bonds, carbon credits, climate fund eligibility \\
\bottomrule
\end{tabularx}

% ── Next Steps ────────────────────────────────────────────────────────────────────
\section*{Getting Started — Next Steps for Your Organisation}

\subsection*{Step 1: Pilot Assessment (Week 1)}
\noindent
Pick one priority neighbourhood or city district (25–100 km²). We'll assess it end-to-end in 2–3 days, producing interactive
risk maps + AI-generated risk narrative. No cost; shows feasibility in your local context.

\subsection*{Step 2: Stakeholder Alignment (Week 2)}
\noindent
Host a 1-hour workshop with planners, infrastructure teams, insurance partners. Review pilot maps. Agree on data governance,
decision thresholds, and success metrics. Secure budget \& sponsorship.

\subsection*{Step 3: Full Deployment (Weeks 3–6)}
\noindent
Deploy on SaaS or on-premises. Train core team (1–2 days). Ingest data for entire city/region. Publish risk layers.

\subsection*{Step 4: Integration \& Action (Months 2–6)}
\noindent
Direct investment decisions, insurance underwriting, and emergency planning based on risk maps. Measure outcomes
(money saved on misdirected infrastructure, lives protected, uninsured risk reduced).

\begin{infobox}
  \textbf{Questions?} Contact the GeoAI team for a confidential 30-min briefing on fit for your organisation,
  pricing, and deployment timeline. No sales pressure — just evidence-based risk intelligence.
\end{infobox}

% ── Footer note ───────────────────────────────────────────────────────────────
\section*{Disclaimer \& Usage Notes}

\begin{tcolorbox}[
  enhanced, arc=6pt,
  colback=primary!5, colframe=primary!30,
  boxrule=0.6pt, left=10pt, right=10pt, top=8pt, bottom=8pt
]
  \small
  \textbf{Document Type:} Business methodology \& ROI brief \quad
  \textbf{Version:} 1.2 (Business-Optimised) \quad
  \textbf{Date:} April 2026 \\[6pt]
  \textbf{Data \& Licencing:} All core data (Open-Meteo, SRTM, Mapillary) are open or commercially licensed. 
  Risk scores are \textit{decision-support tools} — designed for strategic planning and resource allocation. 
  For life-safety applications (e.g., emergency evacuation routing) or detailed engineering design, validate results 
  with qualified domain experts and local hazard data. Insurance underwriting should integrate GeoAI scores with 
  local loss experience and professional underwriting judgment.
  \\[4pt]
  \textcolor{accentdark}{\textbf{Ready to discuss?}} info@geoai.com \quad \textbf{MENA Regional Lead:} [contact details]
\end{tcolorbox}

\end{document}
